%% sections/methodology.tex
\section{Methodology}
\label{sec:methodology}

We design three experimental systems to evaluate different strategies for building a Konkani AMR parser.

\subsection{Base Model}

All three systems share the same backbone: \texttt{BramVanroy/mbart-large-cc25-ft-amr30-en}, a multilingual seq2seq model based on mBART-large \cite{liu2020multilingual} that has been fine-tuned on AMR 3.0 for English AMR parsing. The model produces AMR graphs in a SPRING-style linearisation \cite{bevilacqua2021one} that uses pointer tokens (\texttt{<pointer:N>}) to encode reentrancies and relational brackets (\texttt{<rel>} / \texttt{</rel>}) to delimit argument spans. We decode these into standard Penman notation via a deterministic post-processing step.

We select this model because it is—to our knowledge—the most capable publicly available multilingual AMR parser, trained to parse semantics from a non-English source language using mBART's multilingual representations. We use the Hindi language code (\texttt{hi\_IN}) as a proxy for Konkani in all experiments, as it is the closest Indic language code supported by the mBART-50 tokeniser.

\subsection{System 1: Zero-Shot Baseline}

The first system applies the pre-trained model directly to Konkani sentences without any task-specific adaptation. This baseline measures how much the multilingual representations of a model trained on AMR 3.0 (English) can be exploited for parsing a related but unseen language. No parameters are updated; we run inference on all 1,100 sentences in the dataset with beam search ($k=4$).

\subsection{System 2: Supervised Fine-Tuning on KonkaniAMR}

The second system fine-tunes the base model directly on our KonkaniAMR training split (880 sentence–graph pairs). The target output for each example is the gold Penman string (post-processed from the SPRING linearisation). Fine-tuning is performed with the Hugging Face \texttt{Seq2SeqTrainer}:

\begin{itemize}
    \item \textbf{Epochs:} 20
    \item \textbf{Batch size:} 4 per device, gradient accumulation steps 4 (effective batch size 16)
    \item \textbf{Optimiser:} AdamW, learning rate $5 \times 10^{-5}$, weight decay 0.01
    \item \textbf{Warmup:} 100 steps
    \item \textbf{Max source length:} 128 tokens; \textbf{max target length:} 512 tokens
    \item \textbf{Mixed precision:} \texttt{fp16} on GPU
\end{itemize}

The best checkpoint is selected by minimum validation loss on the 55-sentence validation set. The evaluation uses the 165-sentence test split.

\subsection{System 3: Language-Adaptive Pre-training + Fine-Tuning}

The third system adds a Konkani language-adaptive pre-training stage before AMR fine-tuning. This two-step pipeline is motivated by the hypothesis that exposing the model to a large volume of Konkani text before task-specific training will improve its internal representations of Konkani morphology and vocabulary, leading to better AMR generation.

\textbf{Step 1 — Continued pre-training on Konkani.} We continue pre-training the base model on the 189,332-sentence Konkani corpus. The model is trained with a reconstruction objective: for each sentence, the model is provided the tokenised input and trained to reproduce the same sequence as output, using \texttt{DataCollatorForSeq2Seq}. Training is performed for 5 epochs with:
\begin{itemize}
    \item \textbf{Batch size:} 16 per device, gradient accumulation steps 4 (effective batch size 64)
    \item \textbf{Optimiser:} AdamW, learning rate $3 \times 10^{-5}$, weight decay 0.01
    \item \textbf{Warmup:} 6\% of total steps
    \item \textbf{Mixed precision:} \texttt{fp16} on GPU (A100 80\,GB)
    \item \textbf{Training time:} approximately 1.5 hours
\end{itemize}
The resulting Konkani-adapted model is saved and used as the initialisation for Step 2.

\textbf{Step 2 — AMR fine-tuning.} The Konkani-adapted model is then fine-tuned on KonkaniAMR using the identical hyperparameters and data split as System 2, enabling a direct comparison of the effect of language-adaptive pre-training.

\subsection{Inference and Decoding}

At inference time, all systems use beam search with $k=4$ beams and a forced beginning-of-sequence token set to \texttt{en\_XX} (the language code used during AMR fine-tuning of the base model). The raw output is decoded from the SPRING linearisation format into Penman notation using a deterministic pointer-resolution step, followed by heuristic clean-up to handle common artefacts (literal value brackets, malformed empty nodes, missing spaces before pointers).

\subsection{Evaluation}

All systems are evaluated using Smatch \cite{cai2013smatch}, which computes precision (P), recall (R), and F1 over the set of AMR triples. We report:
\begin{itemize}
    \item \textbf{F1 (parseable):} average Smatch F1 over examples for which both the gold and predicted graphs are valid Penman strings parseable by the \texttt{penman} library.
    \item \textbf{F1 (all):} average Smatch F1 treating all unparseable predictions as F1 = 0, computed over the full test set.
\end{itemize}
This dual reporting reflects both the quality of valid parses and the overall parser reliability. A parser that produces structurally corrupt output on many examples is penalised by the second metric even if its valid outputs score well.
